% Part: first-order-logic
% Chapter: syntax-and-semantics
% Section: formation-sequences

\documentclass[../../../include/open-logic-section]{subfiles}

\begin{document}

\olfileid{fol}{syn}{fseq}

\olsection{Constructiereeksen}

Het definiëren van !!{formula}s met een inductieve definitie en de
aanvullende techniek om eigenschappen van !!{formula}s met inductie te
bewijzen vormen een elegante en efficiënte aanpak. Het kan echter ook
nuttig zijn de constructie van !!{formula}s van onderaf en stap voor
stap te bekijken. Dat doen we hier aan de hand van het begrip
\emph{constructiereeks}.
%
Om te laten zien hoe termen en !!{formula}s op deze manier kunnen worden
ingevoerd zonder naar hun inductieve definities te verwijzen, voeren we
eerst het begrip in van een willekeurige symboolrij uit een taal~$\Lang L$.

\begin{defn}[Symboolrijen]
\ollabel{defn:string}
Stel dat $\Lang L$ een eerste-ordetaal is. Een \emph{$\Lang
L$-symboolrij} is een eindige rij symbolen van~$\Lang L$. Wanneer uit de
context duidelijk is dat de taal~$\Lang L$ vastligt, noemen we een
$\Lang L$-symboolrij kortweg een \emph{symboolrij}.
\end{defn}

\begin{ex}
Voor iedere eerste-ordetaal $\Lang L$ zijn alle
$\Lang L$-!!{formula}s $\Lang L$-symboolrijen, maar de omkering geldt
niet. Zo is \[)(\Obj v_0\lif\lexists\] wel een
$\Lang L$-symboolrij, maar geen $\Lang L$-!!{formula}.
\end{ex}

\begin{defn}[Constructiereeksen voor termen]
\ollabel{defn:fseq-trm}
Een eindige rij $\Lang L$-symboolrijen $\tuple{t_0,\dotsc,t_n}$ is een
\emph{constructiereeks} voor een term $t$ als $t \ident t_n$ en voor
iedere $i \leq n$ geldt: $t_i$ is !!a{variable} of !!a{constant}, of
$\Lang L$ bevat een $k$-plaatsige !!{function}~$f$ en er bestaan
$m_0,\dotsc,m_k < i$ zodanig dat $t_i \ident
f(t_{m_0},\dotsc,t_{m_k})$. Waar het onderscheid nodig is, noemen we
een constructiereeks voor een term een \emph{termconstructiereeks}.
\end{defn}

\begin{ex}
De rij
\[
    \tuple{\Obj c_0, \Obj v_0, \Atom{\Obj f^2_0}{\Obj c_0, \Obj v_0}, \Atom{\Obj f^1_0}{\Atom{\Obj f^2_0}{\Obj c_0, \Obj v_0}}}
\]
is een constructiereeks voor de term $\Atom{\Obj f^1_0}{\Atom{\Obj
f^2_0}{\Obj c_0, \Obj v_0}}$, evenals
\[
    \tuple{\Obj v_0, \Obj c_0, \Atom{\Obj f^2_0}{\Obj c_0, \Obj v_0}, \Atom{\Obj f^1_0}{\Atom{\Obj f^2_0}{\Obj c_0, \Obj v_0}}}.
\]
\end{ex}

\begin{defn}[Constructiereeksen voor formules]
\ollabel{defn:fseq-frm}
Een eindige rij $\Lang L$-symboolrijen $\tuple{!A_0,\dotsc,!A_n}$
is een \emph{constructiereeks} voor~$!A$ als $!A \ident !A_n$ en
voor iedere $i \leq n$ geldt: $!A_i$ is een atomaire !!{formula}, of
er bestaan $j,k < i$ en !!a{variable}~$x$ zodanig dat een van de
volgende gevallen geldt:
\begin{enumerate}
    \tagitem{prvNot}{$!A_i \ident \lnot !A_j$.}{}%
    \tagitem{prvAnd}{$!A_i \ident (!A_j \land !A_k)$.}{}%
    \tagitem{prvOr}{$!A_i \ident (!A_j \lor !A_k)$.}{}%
    \tagitem{prvIf}{$!A_i \ident (!A_j \lif !A_k)$.}{}%
    \tagitem{prvIff}{$!A_i \ident (!A_j \liff !A_k)$.}{}%
    \tagitem{prvAll}{$!A_i \ident \lforall[x][!A_j]$.}{}%
    \tagitem{prvEx}{$!A_i \ident \lexists[x][!A_j]$.}{}%
\end{enumerate}
Waar het onderscheid nodig is, noemen we een constructiereeks voor een
formule een \emph{formuleconstructiereeks}.
\end{defn}

\begin{ex}
\[
\tuple{
    \Atom{\Obj A^1_0}{\Obj v_0},
    \Atom{\Obj A^1_1}{\Obj c_1},
    (\Atom{\Obj A^1_1}{\Obj c_1} \land \Atom{\Obj A^1_0}{\Obj v_0}),
    \lexists[\Obj v_0][(\Atom{\Obj A^1_1}{\Obj c_1} \land \Atom{\Obj A^1_0}{\Obj v_0})]
}
\]
is een constructiereeks voor $\lexists[\Obj v_0][(\Atom{\Obj A^1_1}{\Obj
c_1} \land \Atom{\Obj A^1_0}{\Obj v_0})]$, evenals
\begin{multline*}
\tuple{
    \Atom{\Obj A^1_0}{\Obj v_0},
    \Atom{\Obj A^1_1}{\Obj c_1},
    (\Atom{\Obj A^1_1}{\Obj c_1} \land \Atom{\Obj A^1_0}{\Obj v_0}),
    \Atom{\Obj A^1_1}{\Obj c_1},\\
    \lforall[\Obj v_1][\Atom{\Obj A^1_0}{\Obj v_0}],
    \lexists[\Obj v_0][(\Atom{\Obj A^1_1}{\Obj c_1} \land \Atom{\Obj A^1_0}{\Obj v_0})]
}.
\end{multline*}
%
Zoals uit het tweede voorbeeld blijkt, mogen constructiereeksen
``rommel'' bevatten: !!{formula}s die overtollig zijn of niet aan de
constructie bijdragen.
\end{ex}

\begin{prop}\ollabel{prop:formed}
Iedere !!{formula}~$!A$ in~$\Frm[L]$ heeft een constructiereeks.
\end{prop}

\begin{proof}
Stel dat $!A$ atomair is. Dan is de rij $\tuple{!A}$ een
constructiereeks voor~$!A$.
%
Stel nu dat $!B$ en~$!C$ respectievelijk de constructiereeksen
$\tuple{!B_0,\dotsc,!B_n}$ en $\tuple{!C_0,\dotsc,!C_m}$
hebben.
%
\begin{enumerate}
  \tagitem{prvNot}{Als $!A \ident \lnot !B$, dan is
      $\tuple{!B_0,\dotsc,!B_n,\lnot !B_n}$
      een constructiereeks voor~$!A$.}{}
  \tagitem{prvAnd}{Als $!A \ident (!B \land !C)$, dan is
      $\tuple{!B_0,\dotsc,!B_n,!C_0,\dotsc,!C_m,(!B_n \land !C_m)}$
      een constructiereeks voor~$!A$.}{}
  \tagitem{prvOr}{Als $!A \ident (!B \lor !C)$, dan is
      $\tuple{!B_0,\dotsc,!B_n,!C_0,\dotsc,!C_m,(!B_n \lor !C_m)}$
      een constructiereeks voor~$!A$.}{}
  \tagitem{prvIf}{Als $!A \ident (!B \lif !C)$, dan is
      $\tuple{!B_0,\dotsc,!B_n,!C_0,\dotsc,!C_m,(!B_n \lif !C_m)}$
      een constructiereeks voor~$!A$.}{}
  \tagitem{prvIff}{Als $!A \ident (!B \liff !C)$, dan is
      $\tuple{!B_0,\dotsc,!B_n,!C_0,\dotsc,!C_m,(!B_n \liff !C_m)}$
      een constructiereeks voor~$!A$.}{}
  \tagitem{prvAll}{Als $!A \ident \lforall[x][!B]$, dan is
      $\tuple{!B_0,\dotsc,!B_n,\lforall[x][!B_n]}$
      een constructiereeks voor~$!A$.}{}
  \tagitem{prvEx}{Als $!A \ident \lexists[x][!B]$, dan is
      $\tuple{!B_0,\dotsc,!B_n,\lexists[x][!B_n]}$
      een constructiereeks voor~$!A$.}{}
\end{enumerate}
Volgens het inductieprincipe voor !!{formula}s heeft iedere
!!{formula} een constructiereeks.
\end{proof}

Ook de omgekeerde richting kan worden bewezen. Dat is van belang omdat
hieruit blijkt dat onze twee manieren om formules te definiëren
equivalent zijn: ze leveren dezelfde resultaten op. Bovendien kunnen we
stellingen over formules dan bewijzen met gewone inductie naar de lengte
van constructiereeksen.

\begin{lem}
\ollabel{lem:fseq-init}
Stel dat $\tuple{!A_0,\dotsc,!A_n}$ een constructiereeks voor~$!A_n$
is en dat $k \leq n$. Dan is $\tuple{!A_0,\dotsc,!A_k}$ een
constructiereeks voor~$!A_k$.
\end{lem}

\begin{proof}
Opgave.
\end{proof}

\begin{prob}
Bewijs \olref[fol][syn][fseq]{lem:fseq-init}.
\end{prob}

\begin{thm}
\ollabel{thm:fseq-frm-equiv}
$\Frm[L]$ is de verzameling van alle $\Lang L$-symboolrijen~$!A$
waarvoor een formuleconstructiereeks voor~$!A$ bestaat.
\end{thm}

\begin{proof}
Laat $F$ de verzameling zijn van alle symboolrijen in de taal~$\Lang L$
die een constructiereeks hebben. Uit
\olref[fol][syn][fseq]{prop:formed} weten we dat $\Frm[L] \subseteq F$;
we bewijzen nu de omgekeerde inclusie.

Stel dat $!A$ een constructiereeks $\tuple{!A_0,\dotsc,!A_n}$ heeft.
We bewijzen dat $!A \in \Frm[L]$ met sterke inductie naar~$n$.
Onze inductiehypothese luidt dat iedere symboolrij met een
constructiereeks van lengte $m < n$ tot $\Frm[L]$ behoort.
Volgens de definitie van een constructiereeks is $!A \ident !A_n$
atomair, of er bestaan $j,k < n$ waarvoor een van de volgende gevallen
geldt:
\begin{enumerate}
    \tagitem{prvNot}{$!A \ident \lnot !A_j$.}{}%
    \tagitem{prvAnd}{$!A \ident (!A_j \land !A_k)$.}{}%
    \tagitem{prvOr}{$!A \ident (!A_j \lor !A_k)$.}{}%
    \tagitem{prvIf}{$!A \ident (!A_j \lif !A_k)$.}{}%
    \tagitem{prvIff}{$!A \ident (!A_j \liff !A_k)$.}{}%
    \tagitem{prvAll}{$!A \ident \lforall[x][!A_j]$.}{}%
    \tagitem{prvEx}{$!A \ident \lexists[x][!A_j]$.}{}%
\end{enumerate}
We onderscheiden nu de gevallen. Als $!A$ atomair is, dan is
$!A_n \in \Frm[L_0]$. Stel vervolgens dat
$!A \equiv (!A_j \land !A_k)$. Volgens
\olref[fol][syn][fseq]{lem:fseq-init} zijn
$\tuple{!A_0,\dotsc,!A_j}$ en $\tuple{!A_0,\dotsc,!A_k}$
respectievelijk constructiereeksen voor $!A_j$ en~$!A_k$. Omdat dit
echte beginstukken van de constructiereeks voor~$!A$ zijn, hebben zij
beide lengte kleiner dan~$n$. Volgens de inductiehypothese behoren
$!A_j$ en~$!A_k$ dus tot~$\Frm[L_0]$ en volgens de definitie van
!!a{formula} geldt hetzelfde voor $(!A_j \land !A_k)$. De overige
gevallen verlopen analoog.
\end{proof}

Termconstructiereeksen hebben eigenschappen die vergelijkbaar zijn met
die van constructiereeksen voor !!{formula}s.

\begin{prop}
\ollabel{prop:fseq-trm-equiv}
$\Trm[L]$ is de verzameling van alle $\Lang L$-symboolrijen $t$
waarvoor een termconstructiereeks voor~$t$ bestaat.
\end{prop}

\begin{proof}
Opgave.
\end{proof}

\begin{prob}
Bewijs \olref[fol][syn][fseq]{prop:fseq-trm-equiv}.
Aanwijzing: volg een strategie die vergelijkbaar is met die in het
bewijs van \olref[fol][syn][fseq]{thm:fseq-frm-equiv}.
\end{prob}

In constructiereeksen kunnen twee soorten ``rommel'' voorkomen:
herhaalde elementen en elementen die niet relevant zijn voor de
constructie van de formule of term. Beide soorten kunnen we uitsluiten
door minimale constructiereeksen te beschouwen.

\begin{defn}[Minimale constructiereeksen]
\ollabel{defn:minimal-fseq}
Een constructiereeks $\tuple{!A_0, \ldots, !A_n}$ voor een
formule~$!A$ is een \emph{minimale constructiereeks} voor~$!A$
als voor iedere andere constructiereeks~$s$ voor~$!A$ geldt dat de
lengte van~$s$ groter dan of gelijk aan~$n+1$ is.

Op dezelfde wijze is een constructiereeks $\tuple{t_0, \ldots, t_n}$
voor een term~$t$ een \emph{minimale constructiereeks}
voor~$t$ als voor iedere andere constructiereeks~$s$ voor~$t$ geldt dat
de lengte van~$s$ groter dan of gelijk aan~$n+1$ is.
\end{defn}

Merk op dat een formule of term meer dan één minimale
constructiereeks kan hebben, maar dat deze reeksen precies dezelfde
symboolrijen bevatten.

\begin{prop}
\ollabel{prop:subformula-equivs}
De volgende uitspraken zijn equivalent:
\begin{enumerate}
    \item $!B$ is een sub-!!{formula} van~$!A$.
    \item $!B$ komt voor in iedere constructiereeks van~$!A$.
    \item $!B$ komt voor in een minimale constructiereeks van~$!A$.
\end{enumerate}
\end{prop}

\begin{proof}
Opgave.
\end{proof}

\begin{prob}
Bewijs \olref[fol][syn][fseq]{prop:subformula-equivs}.
\end{prob}

\begin{history}
Constructiereeksen werden door Raymond Smullyan ingevoerd in zijn
leerboek \emph{First-Order Logic} \citep{Smullyan1968}.
Verdere eigenschappen van constructiereeksen werden vastgesteld door
\citet{Zuckerman1973}.
\end{history}
\end{document}
